%% kapa: exam=1 round=21 year=2010 period=1 form=A subject=civil_law subject_ko=민법 seq=20 number=20
%% source: Q-Net 감정평가사 기출문제 (cst00309) artlSeq=5024181 file=extracted/1008/5024181/2010년 제21회 감정평가사 1차시험 기출문제/1교시.hwp
%% conversion: kordoc 4.24.51 -> tools/gen1.py (auto); review=proofread-sonnet; flags=none
\begin{question}{20}
소멸시효와 제척기간에 관한 설명으로 옳지 않은 것은? (다툼이 있으면 판례에 의함)
\begin{choices}
  \item 소멸시효가 완성되면 그 권리는 그 때부터 소멸의 효과가 발생한다.
  \item 소멸시효의 이익은 시효완성 후에 포기할 수 있으나, 제척기간의 경우에는 포기가 인정되지 않는다.
  \item 소멸시효는 법률행위로 이를 단축 또는 경감할 수 있다.
  \item 매매의 일방예약에서의 예약완결권은 예약이 성립한 때로부터 10년 내에 행사하여야 하는데, 이는 제척기간이다.
  \item 제척기간에는 기간의 중단이 있을 수 없다.
\end{choices}
\end{question}
\answer{1}
